% GENERATED FILE. Edit canonical lessons, then run npm run generate:books.
% canonical-source-hash: fnv1a64:48ccd046ea9724d7
% canonical-lessons: TA-C485-kasayam, TA-C485-cukku, TA-C485-kanci, TA-C485-tailam, TA-C485-kalimpu

\chapter{Home Remedies}
\label{ch:ta-home-remedies}
\hlchaptermodality{485}

\begin{chapteropening}
\textbf{By the end of this chapter:} \emph{I can talk about decoctions, dry ginger, gruel, balms and ointments.}

Complete the last lesson of chapter 485: I can talk about decoctions, dry ginger, gruel, balms and ointments.
\end{chapteropening}

\section[\texorpdfstring{\ta{கஷாயம்} (kaṣāyam)}{கஷாயம் (kaṣāyam)}]{\ta{கஷாயம்} (kaṣāyam) — a herbal decoction}

\label{lesson:TA-C485-kasayam}



\begin{quote}
Before the new one: say the Tamil for to fall out, then the Tamil for to turn grey.
\end{quote}

\subsection*{You'll want to know: \ta{கஷாயம்}}
\textbf{\ta{கஷாயம்}} — \emph{kaṣāyam} — \textquotedblleft{}a herbal decoction\textquotedblright{}.

\begin{quote}
\ta{சளிக்குப்} \ta{பாட்டி} \ta{கஷாயம்} \ta{செய்து} \ta{கொடுத்தாள்}.
\end{quote}

\begin{quote}
\emph{caḷikkup pāṭṭi kaṣāyam ceytu koṭuttāḷ.}
\end{quote}

\begin{quote}
Grandmother made a decoction for the cold.
\end{quote}

Written with the grantha \ta{ஷ}, as Tamil writes it today.

\begin{cousinweb}
\textbf{\ta{கஷாயம்}} (\emph{kaṣāyam}) is borrowed from Sanskrit \emph{kaṣāya}, \textquotedblleft{}astringent; an astringent juice\textquotedblright{} (CDIAL 2974).
\end{cousinweb}

\begin{grammarlens}[title={What you've built}]
One more word for describing what you lived and felt.
\end{grammarlens}

\subsection*{Your turn}
\begin{itemize}
\raggedright
  \item \emph{Say it:} \emph{kaṣāyam}
  \item \emph{Say it:} \emph{kaṣāyam}, once more
  \item \emph{Say it:} \emph{caḷikkup pāṭṭi kaṣāyam ceytu koṭuttāḷ.} — Grandmother made a decoction for the cold.
  \item \emph{Recall:} say the Tamil for to fall out, then the Tamil for to turn grey, then say \emph{kaṣāyam} again
\end{itemize}

\begin{tcolorbox}[breakable,colback=teal!4,colframe=teal!35!black,title={Before you move on}]
What does \ta{கஷாயம்} mean? (\textquotedblleft{}A herbal decoction\textquotedblright{}.) Say it once more.
\end{tcolorbox}
\section[\texorpdfstring{\ta{சுக்கு} (cukku)}{சுக்கு (cukku)}]{\ta{சுக்கு} (cukku) — dry ginger}

\label{lesson:TA-C485-cukku}



\begin{quote}
Before the new one: say the Tamil for to turn grey, then the Tamil for a herbal decoction.
\end{quote}

\subsection*{You'll want to know: \ta{சுக்கு}}
\textbf{\ta{சுக்கு}} — \emph{cukku} — \textquotedblleft{}dry ginger\textquotedblright{}.

\begin{quote}
\ta{சுக்கு} \ta{போட்ட} \ta{காபி} \ta{தொண்டைக்கு} \ta{நல்லது}.
\end{quote}

\begin{quote}
\emph{cukku pōṭṭa kāpi toṇṭaikku nallatu.}
\end{quote}

\begin{quote}
Coffee with dry ginger is good for the throat.
\end{quote}

Fresh ginger is \ta{இஞ்சி}; dried, it is \ta{சுக்கு}.

\begin{cousinweb}
The origin of \textbf{\ta{சுக்கு}} (\emph{cukku}) is not settled; it is sometimes compared with Sanskrit \emph{śuṇṭhī}, \textquotedblleft{}dry ginger\textquotedblright{} (CDIAL 12515), but the link is uncertain.
\end{cousinweb}

\begin{grammarlens}[title={What you've built}]
One more word for describing what you lived and felt.
\end{grammarlens}

\subsection*{Your turn}
\begin{itemize}
\raggedright
  \item \emph{Say it:} \emph{cukku}
  \item \emph{Say it:} \emph{cukku}, once more
  \item \emph{Say it:} \emph{cukku pōṭṭa kāpi toṇṭaikku nallatu.} — Coffee with dry ginger is good for the throat.
  \item \emph{Recall:} say the Tamil for to turn grey, then the Tamil for a herbal decoction, then say \emph{cukku} again
\end{itemize}

\begin{tcolorbox}[breakable,colback=teal!4,colframe=teal!35!black,title={Before you move on}]
What does \ta{சுக்கு} mean? (\textquotedblleft{}Dry ginger\textquotedblright{}.) Say it once more.
\end{tcolorbox}
\section[\texorpdfstring{\ta{கஞ்சி} (kañci)}{கஞ்சி (kañci)}]{\ta{கஞ்சி} (kañci) — rice gruel; starch}

\label{lesson:TA-C485-kanci}



\begin{quote}
Before the new one: say the Tamil for a herbal decoction, then the Tamil for dry ginger.
\end{quote}

\subsection*{You'll want to know: \ta{கஞ்சி}}
\textbf{\ta{கஞ்சி}} — \emph{kañci} — \textquotedblleft{}rice gruel; starch\textquotedblright{}.

\begin{quote}
\ta{காய்ச்சலின்போது} \ta{அம்மா} \ta{எனக்குக்} \ta{கஞ்சி} \ta{கொடுத்தாள்}.
\end{quote}

\begin{quote}
\emph{kāyccaliṉpōtu ammā eṉakkuk kañci koṭuttāḷ.}
\end{quote}

\begin{quote}
When I had a fever, mother gave me rice gruel.
\end{quote}

Also starch for clothes: \ta{கஞ்சி} \ta{போட்ட} \ta{சட்டை}, a starched shirt.

\begin{cousinweb}
\textbf{\ta{கஞ்சி}} (\emph{kañci}) is native Dravidian (DEDR 1107), the water poured off rice and the gruel made from it.
\end{cousinweb}

\begin{grammarlens}[title={What you've built}]
One more word for describing what you lived and felt.
\end{grammarlens}

\subsection*{Your turn}
\begin{itemize}
\raggedright
  \item \emph{Say it:} \emph{kañci}
  \item \emph{Say it:} \emph{kañci}, once more
  \item \emph{Say it:} \emph{kāyccaliṉpōtu ammā eṉakkuk kañci koṭuttāḷ.} — When I had a fever, mother gave me rice gruel.
  \item \emph{Recall:} say the Tamil for a herbal decoction, then the Tamil for dry ginger, then say \emph{kañci} again
\end{itemize}

\begin{tcolorbox}[breakable,colback=teal!4,colframe=teal!35!black,title={Before you move on}]
What does \ta{கஞ்சி} mean? (\textquotedblleft{}Rice gruel; starch\textquotedblright{}.) Say it once more.
\end{tcolorbox}
\section[\texorpdfstring{\ta{தைலம்} (tailam)}{தைலம் (tailam)}]{\ta{தைலம்} (tailam) — a medicated oil; a balm}

\label{lesson:TA-C485-tailam}



\begin{quote}
Before the new one: say the Tamil for dry ginger, then the Tamil for rice gruel.
\end{quote}

\subsection*{You'll want to know: \ta{தைலம்}}
\textbf{\ta{தைலம்}} — \emph{tailam} — \textquotedblleft{}a medicated oil; a balm\textquotedblright{}.

\begin{quote}
\ta{தலைவலிக்கு} \ta{அவள்} \ta{நெற்றியில்} \ta{தைலம்} \ta{பூசினாள்}.
\end{quote}

\begin{quote}
\emph{talaivalikku avaḷ neṟṟiyil tailam pūciṉāḷ.}
\end{quote}

\begin{quote}
For her headache she rubbed balm on her forehead.
\end{quote}

\ta{தைலம்} \ta{பூசு}, to rub on balm.

\begin{cousinweb}
\textbf{\ta{தைலம்}} (\emph{tailam}) is borrowed from Sanskrit \emph{taila}, \textquotedblleft{}sesame oil, oil\textquotedblright{} (CDIAL 5958).
\end{cousinweb}

\begin{grammarlens}[title={What you've built}]
One more word for describing what you lived and felt.
\end{grammarlens}

\subsection*{Your turn}
\begin{itemize}
\raggedright
  \item \emph{Say it:} \emph{tailam}
  \item \emph{Say it:} \emph{tailam}, once more
  \item \emph{Say it:} \emph{talaivalikku avaḷ neṟṟiyil tailam pūciṉāḷ.} — For her headache she rubbed balm on her forehead.
  \item \emph{Recall:} say the Tamil for dry ginger, then the Tamil for rice gruel, then say \emph{tailam} again
\end{itemize}

\begin{tcolorbox}[breakable,colback=teal!4,colframe=teal!35!black,title={Before you move on}]
What does \ta{தைலம்} mean? (\textquotedblleft{}A medicated oil; a balm\textquotedblright{}.) Say it once more.
\end{tcolorbox}
\section[\texorpdfstring{\ta{களிம்பு} (kaḷimpu)}{களிம்பு (kaḷimpu)}]{\ta{களிம்பு} (kaḷimpu) — an ointment}

\label{lesson:TA-C485-kalimpu}



\begin{quote}
Before the new one: say the Tamil for rice gruel, then the Tamil for a medicated oil.
\end{quote}

\subsection*{You'll want to know: \ta{களிம்பு}}
\textbf{\ta{களிம்பு}} — \emph{kaḷimpu} — \textquotedblleft{}an ointment\textquotedblright{}.

\begin{quote}
\ta{மருத்துவர்} \ta{புண்ணுக்கு} \ta{ஒரு} \ta{களிம்பு} \ta{கொடுத்தார்}.
\end{quote}

\begin{quote}
\emph{maruttuvar puṇṇukku oru kaḷimpu koṭuttār.}
\end{quote}

\begin{quote}
The doctor gave an ointment for the sore.
\end{quote}

Its older sense is verdigris, the green rust on copper.

\begin{cousinweb}
\textbf{\ta{களிம்பு}} (\emph{kaḷimpu}) is native Dravidian (DEDR 1586), where it means verdigris and rust; the sense of ointment is the modern one.
\end{cousinweb}

\begin{grammarlens}[title={What you've built}]
One more word for describing what you lived and felt.
\end{grammarlens}

\subsection*{Your turn}
\begin{itemize}
\raggedright
  \item \emph{Say it:} \emph{kaḷimpu}
  \item \emph{Say it:} \emph{kaḷimpu}, once more
  \item \emph{Say it:} \emph{maruttuvar puṇṇukku oru kaḷimpu koṭuttār.} — The doctor gave an ointment for the sore.
  \item \emph{Recall:} say the Tamil for rice gruel, then the Tamil for a medicated oil, then say \emph{kaḷimpu} again
\end{itemize}

\begin{tcolorbox}[breakable,colback=teal!4,colframe=teal!35!black,title={Before you move on}]
What does \ta{களிம்பு} mean? (\textquotedblleft{}An ointment\textquotedblright{}.) Say it once more.
\end{tcolorbox}
